\chapter{Instructions and Marking Scheme}\label{ch:instructions}

This part simulates the departmental written examination. It has three sections, in the
order in which they are usually set.

\begin{table}[ht]\centering\small
\begin{tabular}{@{}llp{0.55\linewidth}@{}}\toprule
Section & Chapter & Content\\\midrule
A.\ Concepts & \ref{ch:concepts} & short written answers: definitions, complexities, small
  proofs and hand simulations\\
B.\ Tracing & \ref{ch:tracing} & read a program and state its exact output; most items
  hide a twist\\
C.\ Problem solving & \ref{ch:problems} & write code against a specification, with test
  cases, permitted operations and a required time complexity\\\bottomrule
\end{tabular}
\end{table}

\section{How to use this part}
\begin{enumerate}[(i)]
\item \emph{Attempt each item on paper first.} Tracing is a paper skill; reading the
  answer teaches little.
\item \emph{Time yourself.} A realistic budget is 2--3 minutes per concept item, 5--8
  minutes per tracing item and 20--30 minutes per easy-to-medium problem.
\item \emph{Check against Part~III}, which gives the real program output, the reasoning,
  and for problems an idea, a complete commented solution, its complexity and the usual
  mistakes.
\item \emph{Run the code.} Every program in the book was compiled and run; retyping the
  solutions is excellent practice, and altering the inputs shows what each line does.
\end{enumerate}

\section{Conventions used in the problem statements}
\begin{itemize}
\item \textbf{Permitted operations.} Each problem lists the operations you may call.
  Anything else must be written by you, from the representation, if the statement allows
  it. Calling an operation that is not listed (for instance indexing the array inside a
  stack when only push and pop are permitted) forfeits the credit even if the output is
  correct.
\item \textbf{Starter code and constraints.} The box under the statement is the function
  (or the few functions) you have to write; the types and helper operations it uses are
  those of the cited listing. The constraints give the sizes the tests may reach.
\item \textbf{Required.} The time (and sometimes space) complexity needed for full credit.
  A slower correct solution earns partial credit.
\item \textbf{Bonus.} Extra credit for beating the required bound, or for an additional
  property (in place, single pass, $O(1)$ extra space).
\item \textbf{Tests.} Sample inputs and expected outputs. They are illustrations, not the
  full test set: a real judge adds the empty case, the one-element case, duplicates, the
  maximum size, a full fixed-capacity structure, indices that wrap around, and adversarial
  orderings.
\item $n$ is the number of items. Positions are 0-based unless the statement says
  otherwise.
\item Unless stated otherwise, you may assume the memory returned by \cc{malloc}/\cc{new}
  is available, and you must not leak nodes that your code removes.
\end{itemize}

\section{A sample marking scheme}
The departmental scheme differs from term to term; the following is a plausible shape,
useful for self-assessment.

\begin{table}[ht]\centering\small
\begin{tabular}{@{}lr@{}}\toprule
Criterion & Share\\\midrule
Correct on all tests, including edge cases & 50\%\\
Meets the required complexity & 20\%\\
Uses only permitted operations (abstraction respected) & 10\%\\
No memory errors, no leaks; \cc{head}/\cc{tail}/\cc{size}/\cc{front} consistent & 10\%\\
Readability and comments & 10\%\\
Bonus for a faster or leaner solution & up to $+20\%$ of the item\\\bottomrule
\end{tabular}
\end{table}

\para{Exam technique}
\begin{enumerate}[(i)]
\item Write the invariant or the picture before the loop. A node-and-arrow sketch of the
  three cases (empty, one node, general), or an array with \cc{front} and \cc{size}
  marked, catches most bugs.
\item Handle the edge cases explicitly, with their own \cc{if}: empty structure, one
  element, first, last, not found, duplicates, everything equal, full, wrap-around.
\item State the complexity in a comment. If you know a faster approach but cannot finish
  it, submit the correct slower one and say so.
\item Never call a library routine that does the very thing being asked (e.g.\
  \cc{std::reverse} in a problem about reversing); do not depend on headers beyond
  \cc{stdio}/\cc{stdlib}/\cc{iostream} unless the statement permits.
\item For a stack, queue or deque problem, first ask \emph{which end} each item leaves
  from. ``The most recent'' means a stack; ``the oldest'' means a queue; ``both'', or
  ``the oldest and also the most recent'', means a deque.
\end{enumerate}
